% !TeX root = ../../Book2.tex
% 纲要标题：### 4.5 Jordan 理论
% 纲要位置：计划纲要.md，第 942 行。

\section{Jordan 理论}
\label{b2:sec:0405}

有理标准形把算子分成循环部分，并在每一部分保留一个关系多项式。如果这些多项式在基域上都分解为一次因子，就可以再按各个根拆开，并用连续施加 \(T-\lambda I\) 得到的向量选基。这样得到的矩阵只有特征值和少量的一，其块大小还可以从核的维数直接读出。本节始终令 \(V\) 为域 \(k\) 上的非零有限维空间，\(T\in\operatorname{End}_k(V)\)，并记 \(n=\dim_kV\)。零空间对应空的块分解，可单独并入结论。

% 纲要内容（仅作写作计划，尚非教材正文）：
%
% * 分裂域上的初等因子
% * Jordan 块及其唯一性
% * 基域扩张如何改变标准形
% * 不变因子在扩域下的保持与相似关系的下降
%
% ---
%

\subsection{分裂域上的初等因子}
\label{b2:subsec:0405:01}

由\cref{b2:thm:rational-canonical-form,b2:prop:operator-characteristic-minimal}，算子模 \(V_T\) 的非平凡不变因子可以写成
\[
f_1\mid f_2\mid\cdots\mid f_s,\qquad
V_T\cong\bigoplus_{j=1}^s k[t]/(f_j),
\]
其中各 \(f_j\) 首一且次数为正，并有
\[
\chi_T=\prod_{j=1}^s f_j,\qquad \mu_T=f_s.
\]
因为每个 \(f_j\) 都整除 \(f_s\)，\(\chi_T\) 与 \(\mu_T\) 具有完全相同的首一不可约因子。因此它们是否在 \(k\) 上分裂为一次因子的乘积，是同一个条件；这里不要求这些一次因子互不相同。

\ProseParagraphBreak
不同特征值的一次因子幂两两互素，所以第一章的 Bézout 核分解\cref{b2:prop:coprime-operator-decomposition}已能把它们对应的广义特征子空间分开。现在应用主理想整环模的分类，是要进一步确定每个主空间内部有多少个循环部分，以及各部分的长度。

\begin{proposition}\label{b2:prop:jordan-splitting-primary}
设 \(V\ne0\) 是域 \(k\) 上的有限维向量空间，\(T\in\operatorname{End}_k(V)\)。假设 \(T\) 的最小多项式 \(\mu_T\) 在 \(k\) 上分裂，写成
\[
\mu_T(t)=\prod_{\lambda\in\Lambda}(t-\lambda)^{a_\lambda},
\qquad a_\lambda\ge1,
\]
其中 \(\Lambda\subseteq k\) 为有限集，所列 \(\lambda\) 互不相同，则
\[
V=\bigoplus_{\lambda\in\Lambda}G_\lambda(T),
\qquad G_\lambda(T)=\ker\bigl((T-\lambda I_V)^{a_\lambda}\bigr),
\]
而每个广义特征子空间 \(G_\lambda(T)\) 作为 \(k[t]\) 模可以分解为
\[
G_\lambda(T)\cong
\bigoplus_{j=1}^{r_\lambda}k[t]/\bigl((t-\lambda)^{d_{\lambda,j}}\bigr),
\qquad d_{\lambda,j}\ge1,
\qquad \max_j d_{\lambda,j}=a_\lambda.
\]
这些幂正是 \(V_T\) 的初等因子。子空间 \(G_\lambda(T)\) 由 \(T\) 唯一确定；其中循环子模的具体选取未必唯一，但各个指数及其出现次数唯一。
\end{proposition}
\begin{proof}
将这些互素因子幂代入\cref{b2:cor:generalized-eigenspace-decomposition}，直接得到 \(V=\bigoplus_\lambda G_\lambda(T)\)，且 \(G_\lambda(T)=\ker((T-\lambda I_V)^{a_\lambda})\)。因此这些主空间已经是 \(V\) 中由算子确定的子空间，而不只是某个抽象模同构选出的分量。每个主空间都非零，否则删去对应的因子幂后仍能零化全空间，与 \(\mu_T\) 的最小性矛盾。

每个 \(G_\lambda(T)\) 有限维，其算子模有限生成，并被 \((t-\lambda)^{a_\lambda}\) 零化。应用\cref{b2:thm:pid-elementary-divisors}，它的初等因子只能是 \((t-\lambda)^{d_{\lambda,j}}\)，其中 \(1\le d_{\lambda,j}\le a_\lambda\)，并且这些指数及出现次数唯一。一个多项式零化整个 \(V_T\)，当且仅当它零化每个循环商；后者等价于每个初等因子都整除这个多项式。因此全模的最小多项式是这些初等因子的最小公倍式，\(t-\lambda\) 的指数为 \(\max_jd_{\lambda,j}\)。与题设的 \(\mu_T\) 比较，得到这个最大值恰为 \(a_\lambda\)，同时证明所列幂就是 \(V_T\) 的全部初等因子。
\end{proof}

在 \(G_\lambda(T)\) 上，\(N_\lambda=(T-\lambda I_V)|_{G_\lambda(T)}\) 为幂零算子。普通特征子空间只是其中的 \(\ker(T-\lambda I_V)\)；整个主空间还包含需要连续施加多次 \(T-\lambda I_V\) 才变为零的向量。每个循环因子的幂次数记录了这种连续作用的长度，所以分类时不能只记录特征值。

\ProseParagraphBreak
例如，一个算子的初等因子若为
\[
(t-1)^3,\qquad (t-1),\qquad (t+1)^2
\]
且 \(\operatorname{char}k\ne2\)，那么 \(G_1(T)\) 的维数为四，\(G_{-1}(T)\) 的维数为二。前一个空间内有长度不同的两个循环因子。下一小节将说明长度三与长度一怎样在矩阵中表现出来。特征不等于二的条件用来保证 \(1\) 与 \(-1\) 是两个不同的根。

\subsection{Jordan 块及其唯一性}
\label{b2:subsec:0405:02}

回顾第一章采用的 \term[Jordan-kuai]{Jordan 块}[Jordan block]记号：\(J_d(\lambda)\) 的主对角元为 \(\lambda\)，紧贴主对角线的上次对角元均为一，其余元素为零；特别地，\(J_1(\lambda)=(\lambda)\)。例如
\[
J_3(\lambda)=
\begin{pmatrix}
\lambda&1&0\\
0&\lambda&1\\
0&0&\lambda
\end{pmatrix}.
\]
块的方向取决于链向量的排列次序。令 \(N=T-\lambda I_V\)，若 \(v\) 满足 \(N^dv=0\)、\(N^{d-1}v\ne0\)，从链顶 \(v\) 开始施加 \(N\)，得到
\[
v\xrightarrow{\,N\,}Nv\xrightarrow{\,N\,}N^2v
\xrightarrow{\,N\,}\cdots\xrightarrow{\,N\,}N^{d-1}v
\xrightarrow{\,N\,}0.
\]
链只列到最后一个非零向量 \(N^{d-1}v\)；\(d=1\) 时就是 \(v\xrightarrow{N}0\)，\(d=2\) 时就是 \(v\xrightarrow{N}Nv\xrightarrow{N}0\)。写矩阵时，基却按 \(N^{d-1}v,\ldots,Nv,v\) 的逆序排列。这样第一个基向量被 \(N\) 送到零，而 \(N\) 把第 \(j\) 个基向量送到第 \(j-1\) 个（\(2\le j\le d\)）；矩阵中第 \(j\) 列的 \(1\) 就位于第 \(j-1\) 行，即上次对角线。下面的引理验证这些向量确实构成循环子空间的基。

\begin{lemma}\label{b2:lem:jordan-chain-block}
设 \(k\) 是域，\(\lambda\in k\)，\(d\ge1\)。在 \(k[t]/((t-\lambda)^d)\) 上考虑乘以 \(t\) 的算子。按下列顺序排列的陪集是一组 \(k\) 基：
\[
e_1=\overline{(t-\lambda)^{d-1}},\quad
e_2=\overline{(t-\lambda)^{d-2}},\quad\ldots,\quad
e_d=\overline1.
\]
算子在这组基下的矩阵为 Jordan 块 \(J_d(\lambda)\)，其对角元均为 \(\lambda\)，紧邻对角线的上方元素均为 \(1\)，其余元素为 \(0\)。等价地，若 \(V\) 是 \(k\) 向量空间，\(T\in\operatorname{End}_k(V)\)，且循环子空间由 \(v\in V\) 生成并满足 \(\mu_{T,v}=(t-\lambda)^d\)，则其相应基为
\[
(T-\lambda I)^{d-1}v,\ (T-\lambda I)^{d-2}v,\ \ldots,\ v.
\]
\end{lemma}
\begin{proof}
把 \(t\) 换成 \((t-\lambda)+\lambda\)，可知次数小于 \(d\) 的多项式都能唯一地写成 \(1,t-\lambda,\ldots,(t-\lambda)^{d-1}\) 的线性组合。带余除法又说明它们的陪集生成所写商空间。若这些陪集有非平凡线性关系，就有一个次数小于 \(d\) 的非零多项式被 \((t-\lambda)^d\) 整除，矛盾。因此逆序排列后仍是一组基。

令 \(N\) 表示乘以 \(t-\lambda\)。在这组基上，\(Ne_1=0\)，而 \(Ne_j=e_{j-1}\) 对 \(2\le j\le d\) 成立。乘 \(t\) 等于 \(\lambda I+N\)，故其第 \(j\) 列恰为 \(\lambda e_j+e_{j-1}\) 的坐标，得到所述矩阵。

对于等价表述，记 \(W=k[T]v\)。次数 \(d\) 的关系使 \(W\) 由 \(v,Tv,\ldots,T^{d-1}v\) 生成，所以 \(W\) 有限维。将\cref{b2:prop:cyclic-polynomial-quotient}用于 \(T|_W\)，得到 \(k[t]/((t-\lambda)^d)\cong W\)。这个同构把 \(\overline1\) 送到 \(v\)，并保持多项式作用，因而把上述基送到第二组向量。
\end{proof}

这样的向量排列称为一条 \term[Jordan-lian]{Jordan 链}[Jordan chain]：它满足 \((T-\lambda I)e_1=0\) 与 \((T-\lambda I)e_j=e_{j-1}\)。例如在 \(k^3\) 中取
\[
A=\begin{pmatrix}
\lambda&1&1\\
0&\lambda&1\\
0&0&\lambda
\end{pmatrix},\qquad N=A-\lambda I.
\]
对标准基向量有 \(Ne_3=e_1+e_2\)、\(N^2e_3=e_1\) 和 \(N^3e_3=0\)。因此 \(e_1,e_1+e_2,e_3\) 是一组 Jordan 链基；三者线性无关，且 \(A\) 在这组基下成为 \(J_3(\lambda)\)。求这样的基，要求实际找到链顶的向量，再逐次施加 \(N\)。

\begin{theorem}[Jordan 标准形]\label{b2:thm:jordan-canonical-form}
设 \(V\) 是域 \(k\) 上的有限维向量空间，\(T\in\operatorname{End}_k(V)\)。\(T\) 在某组基下的矩阵是 Jordan 块的直和，当且仅当其特征多项式在 \(k\) 上分裂；这也等价于最小多项式在 \(k\) 上分裂。存在这种表示时，各特征值的 Jordan 块大小及其重数唯一，只有块的排列次序可以改变。
\end{theorem}
\begin{proof}
\(V=0\) 时，特征多项式和最小多项式均为 \(1\)，对应空的块分解，结论成立。以下设 \(V\ne0\)。若已有 Jordan 块表示，每个块的特征多项式为 \((t-\lambda)^d\)，故全部块的特征多项式相乘后在 \(k\) 上分裂。反过来，分裂条件使\cref{b2:prop:jordan-splitting-primary}适用；在每个循环因子上取\cref{b2:lem:jordan-chain-block}中的基，再把它们合在一起，便得到整个空间的 Jordan 基。

为证明唯一性，固定 \(\lambda\in k\)，令
\[
c_\lambda(r)=\dim_k\ker\bigl((T-\lambda I_V)^r\bigr)\quad(r\ge1),
\qquad c_\lambda(0)=0.
\]
一个属于 \(\lambda\)、大小为 \(d\) 的块，在 \((T-\lambda I)^r\) 的核中贡献 \(\min(r,d)\) 维，因为链的前 \(\min(r,d)\) 个向量恰被零化。对于属于 \(\mu\ne\lambda\) 的块，\(T-\lambda I=(\mu-\lambda)I+N\)，其中 \(N^d=0\)。记 \(c=\mu-\lambda\ne0\)，有限和
\[
(cI+N)^{-1}=c^{-1}\sum_{j=0}^{d-1}(-c^{-1}N)^j
\]
确为逆算子，所以这样的块对核维数没有贡献。若属于 \(\lambda\) 的块大小为 \(d_{\lambda,1},\ldots,d_{\lambda,r_\lambda}\)，就有
\begin{equation}\label{b2:eq:jordan-kernel-dimensions}
c_\lambda(r)=\sum_{j=1}^{r_\lambda}\min(r,d_{\lambda,j}).
\end{equation}
将每个块画成高度为 \(d_{\lambda,j}\) 的一列格子，\(c_\lambda(r)\) 就是在高度 \(r\) 截平各列后留下的格子总数。由 \(r-1\) 增加到 \(r\)，每根尚未结束的列恰增加一个格子，所以新增层的格子数等于大小至少为 \(r\) 的块数。这正是\cref{b2:fig:pid-primary-layers}中的主分量层数计数，在 \(p=t-\lambda\) 时的解释。更具体地，记 \(N_\lambda=(T-\lambda I_V)|_{G_\lambda(T)}\)，由有限维秩与零度公式，核维数的增量也等于相邻像空间的层商维数：
\[
\begin{aligned}
c_\lambda(r)-c_\lambda(r-1)
&=\dim_k\left(N_\lambda^{r-1}G_\lambda(T)/N_\lambda^rG_\lambda(T)\right)\\
&=\#\{\,j\mid d_{\lambda,j}\ge r\,\}.
\end{aligned}
\]
再取相邻两项之差，便得到大小恰为 \(r\) 的块数。因此全部块大小及其重数都由 \(T\) 的这些核维数决定，不依赖所选基，证明唯一性。
\end{proof}

Jordan 基本身通常不唯一；唯一的是块的分类数据。不同教材若把一放在下次对角，只需在每条链中逆序排列基向量，即得到这里的约定。

本章的分类数据可以这样接起来。对不变因子链 \(f_1\mid\cdots\mid f_s\)，逐个将 \(f_i\) 分解为首一不可约多项式的幂，设 \(f_i=\prod_p p^{e_{p,i}}\)，只记录指数 \(e_{p,i}>0\) 的项。每一项便给出一个初等因子；在分裂情形下，每个 \(p=t-\lambda\)，相应的循环商由上面的引理给出一个 Jordan 块：
\[
\begin{aligned}
f_i=\prod_p p^{e_{p,i}}
&\quad\longmapsto\quad p^{e_{p,i}}\quad(e_{p,i}>0),\\
p^{e_{p,i}}=(t-\lambda)^{e_{p,i}}
&\quad\longmapsto\quad J_{e_{p,i}}(\lambda).
\end{aligned}
\]
原有的重复因子要逐个保留；特征多项式是 \(\prod_i f_i\)，最小多项式是最后一个不变因子 \(f_s\)。

\begin{proposition}\label{b2:prop:jordan-kernel-dimensions}
设 \(V\) 是域 \(k\) 上的有限维向量空间，\(T\in\operatorname{End}_k(V)\) 的特征多项式在 \(k\) 上分裂，\(\lambda\in k\) 是其特征值。记 \(c_\lambda(0)=0\)，\(c_\lambda(r)=\dim_k\ker\bigl((T-\lambda I_V)^r\bigr)\)（\(r\ge1\)）。此时属于 \(\lambda\)、大小恰为 \(r\) 的 Jordan 块数是
\[
2c_\lambda(r)-c_\lambda(r-1)-c_\lambda(r+1).
\]
块的总数为 \(c_\lambda(1)\)，各块大小之和为 \(\dim_k G_\lambda(T)\)，最大块大小为 \(\mu_T\) 中 \(t-\lambda\) 的指数。特别地，\(\lambda\) 的几何重数为块的总数，代数重数为各块大小之和。
\end{proposition}
\begin{proof}
令 \(a_r=c_\lambda(r)-c_\lambda(r-1)\)。由定理证明中的计数，\(a_r\) 是大小至少为 \(r\) 的块数；其中还能进入第 \(r+1\) 层的有 \(a_{r+1}\) 个，所以大小恰为 \(r\) 的块数为
\[
\begin{aligned}
a_r-a_{r+1}
&=\bigl[c_\lambda(r)-c_\lambda(r-1)\bigr]
-\bigl[c_\lambda(r+1)-c_\lambda(r)\bigr]\\
&=2c_\lambda(r)-c_\lambda(r-1)-c_\lambda(r+1).
\end{aligned}
\]
取 \(r=1\)，每个块贡献一个特征向量，故块总数为 \(c_\lambda(1)\)。全部循环商的维数相加给出 \(\dim_k G_\lambda(T)\)，其中最大的指数则由\cref{b2:prop:jordan-splitting-primary}确定。最后，每个块在特征多项式中贡献 \((t-\lambda)^d\)，所以块大小之和也是代数重数。
\end{proof}

例如，一个七维算子只有特征值 \(\lambda\)，已算出
\[
c_\lambda(0),c_\lambda(1),\ldots,c_\lambda(5)
=0,3,5,6,7,7.
\]
\cref{b2:tab:jordan-kernel-differences}把相邻差及其含义列在一起。
\begin{table}[htbp]
\centering
\caption[核维数与 Jordan 块]{核维数及其相邻差分}
\label{b2:tab:jordan-kernel-differences}
\begin{tabular}{@{}ccc l@{}}
\toprule
\(r\) & \(c_\lambda(r)\) & \(c_\lambda(r)-c_\lambda(r-1)\) & 含义 \\
\midrule
1 & 3 & 3 & 三个块 \\
2 & 5 & 2 & 两个块大小至少为二 \\
3 & 6 & 1 & 一个块大小至少为三 \\
4 & 7 & 1 & 一个块大小至少为四 \\
5 & 7 & 0 & 没有块大小至少为五 \\
\bottomrule
\end{tabular}
\end{table}
相邻差的变化说明有一个大小为四的块、一个大小为二的块和一个大小为一的块，没有其他块。它的 Jordan 标准形为
\[
J_4(\lambda)\oplus J_2(\lambda)\oplus J_1(\lambda).
\]
这里无需先求出全部 Jordan 链，就能由核维数确定标准形；若还要写出换基矩阵，则必须进一步求出链向量。

\begin{corollary}\label{b2:cor:jordan-diagonalizable}
设 \(V\) 是域 \(k\) 上的有限维向量空间，\(T\in\operatorname{End}_k(V)\)，则 \(T\) 在 \(k\) 上可对角化，当且仅当其最小多项式 \(\mu_T\) 是 \(k[t]\) 中互不相同的一次因子的乘积；\(V=0\) 时将空乘积视为 \(1\)。
\end{corollary}
\begin{proof}
若 \(\mu_T\) 具有这种分解，Jordan 标准形存在，且由\cref{b2:prop:jordan-kernel-dimensions}，每个特征值的最大块大小均为一，故整矩阵为对角矩阵。反过来，一个对角矩阵的 Jordan 块全为大小一，因而最小多项式的各一次因子只出现一次。
\end{proof}

\subsection{基域扩张与标准形}
\label{b2:subsec:0405:03}

一个在 \(k\) 上不可约的多项式，到了较大的域 \(L\) 中可能分解。要比较相应的算子，应先说明改变了哪个空间：给定 \(A\in M_n(k)\)，把同一矩阵的条目看成 \(L\) 中的元素，使它作用在 \(L^n\) 上，所得算子记为 \(T_L\)。它并不是把原来的 \(k^n\) 直接当作 \(L\) 向量空间。

\ProseParagraphBreak
这种比较不依赖最初选取的 \(k\) 基。若 \(A'=P^{-1}AP\) 且 \(P\in\operatorname{GL}_n(k)\)，同一个 \(P\) 也在 \(\operatorname{GL}_n(L)\) 中，因而 \(A,A'\) 在 \(L\) 上仍相似。下面用矩阵说明标准形怎样变化，无须预先使用向量空间的张量积。

\begin{proposition}\label{b2:prop:scalar-extension-invariant-factors}
设 \(k\subseteq L\) 为域扩张，\(n\ge1\)，\(A\in M_n(k)\)，其在 \(k[t]\) 中的首一不变因子为 \(f_1,\ldots,f_s\)。把 \(A\) 看作 \(L\) 上的矩阵后，首一不变因子仍为这些多项式，只是把系数看成 \(L\) 中的元素。因此最小多项式和特征多项式也保持不变。
\end{proposition}
\begin{proof}
由\cref{b2:prop:operator-polynomial-matrix}，\(tI_n-A\) 是相应算子模的一个展示矩阵。因为 \(k[t]\) 是主理想整环，\cref{b2:thm:smith-normal-form}给出它的 Smith 分解，写成
\[
P(t)(tI_n-A)Q(t)
=\operatorname{diag}(1,\ldots,1,f_1,\ldots,f_s),
\]
其中 \(P,Q\) 是 \(k[t]\) 上的可逆矩阵。它们的逆矩阵也有 \(k[t]\) 中的条目，所以放到 \(L[t]\) 上仍可逆，所写等式原样成立。原有的整除关系 \(f_1\mid\cdots\mid f_s\) 保持成立；每个 \(f_j\) 仍首一且次数为正，因而不会在 \(L[t]\) 中变成单位。这仍是 Smith 标准形。

将\cref{b2:prop:operator-polynomial-matrix}用于 \(L^n\)，并用\cref{b2:thm:pid-elementary-divisors}中不变因子的唯一性，便知同一列 \(f_j\) 就是 \(T_L\) 的不变因子。再由\cref{b2:prop:operator-characteristic-minimal}中的公式 \(\mu_T=f_s\)、\(\chi_T=\prod_j f_j\)，得到最后的断言。
\end{proof}

\begin{corollary}\label{b2:cor:similarity-descent}
设 \(k\subseteq L\) 为域扩张，\(n\ge1\)，\(A,B\in M_n(k)\)。以 \(\sim_k\)、\(\sim_L\) 分别表示在 \(k\)、\(L\) 上的矩阵相似关系，则
\[
A\sim_LB\quad\Longleftrightarrow\quad A\sim_kB.
\]
\end{corollary}
\begin{proof}
若 \(B=P^{-1}AP\)，其中 \(P\in\operatorname{GL}_n(k)\)，则 \(P\) 的逆矩阵仍有 \(k\) 中的条目，故 \(P\) 也属于 \(\operatorname{GL}_n(L)\)，原等式在 \(L\) 中仍成立。因此 \(A\sim_kB\) 蕴含 \(A\sim_LB\)。

反过来，若 \(A\sim_LB\)，由\cref{b2:thm:rational-canonical-form}，它们在 \(L[t]\) 中具有相同的首一不变因子。\cref{b2:prop:scalar-extension-invariant-factors}说明这两列分别就是 \(A,B\) 原来在 \(k[t]\) 中的不变因子。由于 \(k\hookrightarrow L\) 单射，两个 \(k\) 系数多项式在 \(L[t]\) 中相等，也就在 \(k[t]\) 中相等。再次应用有理标准形的相似判据，得到 \(A\sim_kB\)。
\end{proof}

扩域可以使原来写不出的特征值出现，却不会把两个不同的原域相似类合并。例如，两个实矩阵若在复数域上相似，它们已经在实数域上相似。

\ProseParagraphBreak
不过，“保持不变”指不变因子、最小多项式与特征多项式作为系数扩张后的多项式不变，并不指它们的不可约分解不变；初等因子要按新的不可约分解重新组织，分裂后才能据此读出 Jordan 块。例如在 \(\mathbb R\) 上，\((t^2+1)^2\) 是一个不可约多项式的二次幂；到 \(\mathbb C\) 上，它则分解为 \((t-i)^2(t+i)^2\)。多项式还是原来那一个，但初等因子已经改变。

\begin{proposition}\label{b2:prop:scalar-extension-elementary-divisors}
设 \(k\subseteq L\) 为域扩张，\(n\ge1\)，\(A\in M_n(k)\)，\(q^e\) 是 \(A\) 在 \(k[t]\) 中的一个初等因子，其中 \(q\in k[t]\) 首一不可约，\(e\ge1\)。若在 \(L[t]\) 中有分解
\[
q=\prod_{j=1}^u q_j^{h_j},
\]
其中 \(q_j\) 为互不相同的首一不可约多项式，\(h_j\ge1\)，则这个循环因子在 \(L\) 上贡献的初等因子为
\[
q_1^{eh_1},\ldots,q_u^{eh_u}.
\]
若 \(q\) 可分，即在代数闭包中没有重根，则任意扩域后的上述分解都有 \(h_j=1\)，原幂指数 \(e\) 保持不变。若 \(q\) 不可分，则不能保证这一点；一旦某个 \(h_j>1\)，新因子的指数就是 \(eh_j\)。
对原来的每个初等因子作这项操作，就得到 \(A\) 在 \(L\) 上的全部初等因子，原有的重复次数也分别保留。特别地，若 \(L\) 是 \(\chi_A\) 的分裂域，这些因子全部为一次多项式的幂，相应的幂指数就是 Jordan 块大小。
\end{proposition}
\begin{proof}
由\cref{b2:thm:pid-elementary-divisors}，先把算子模分成初等因子对应的循环商，再在每个 \(k[t]/(q^e)\) 中取 \(1,t,\ldots,t^{\deg(q^e)-1}\) 的陪集为基，就把 \(A\) 写成各个友矩阵 \(C(q^e)\) 的直和。同一个换基矩阵在 \(L\) 上仍可逆，而 \(C(q^e)\) 在 \(L\) 上仍表示 \(L[t]/(q^e)\) 中乘 \(t\) 的算子。因此只需分解这个循环商。

多项式 \(q_j^{eh_j}\) 两两互素。对逐项取余的同态
\[
L[t]/(q^e)\longrightarrow
\bigoplus_{j=1}^u L[t]/(q_j^{eh_j}),
\]
核为零，因为能被各个互素因子整除就能被它们的乘积 \(q^e\) 整除。为证明满射，对每个 \(j\) 令 \(Q_j=q^e/q_j^{eh_j}\)，由 Bézout 等式取 \(u_j,v_j\)，使
\[
u_jQ_j+v_jq_j^{eh_j}=1.
\]
给定各坐标的代表多项式 \(g_j\)，多项式 \(\sum_j g_ju_jQ_j\) 在第 \(j\) 项模 \(q_j^{eh_j}\) 后恰为 \(g_j\)，所以给出原像。于是所写同态为保持 \(t\) 作用的同构，得到结论。不同原不可约因子在 \(k[t]\) 中已有 Bézout 等式，在 \(L[t]\) 中仍互素，因此不会在扩域后混入同一项因子；最后再用初等因子的唯一性即可。

若 \(q\) 可分，则 \(q\) 与其形式导数 \(q'\) 在 \(k[t]\) 中互素，其 Bézout 等式放到 \(L[t]\) 中仍成立。因此 \(q\) 在 \(L[t]\) 中也不含重复的不可约因子，否则这个重复因子会同时整除 \(q\) 与 \(q'\)，与互素性矛盾。故每个 \(h_j=1\)，而一般情形中的指数由 \(q^e=\prod_jq_j^{eh_j}\) 得到。
\end{proof}

例如，令实矩阵 \(A=C((t^2+1)^2)\)。它只有一个不变因子 \((t^2+1)^2\)，最小多项式和特征多项式均等于这个四次多项式。在 \(\mathbb C\) 上，不变因子仍只有这一个，但初等因子变为 \((t-i)^2\) 与 \((t+i)^2\)。因此复数域上的 Jordan 标准形为
\[
J_2(i)\oplus J_2(-i).
\]
扩大基域使特征值可以写出来，却没有使这个算子可对角化，因为两个块都仍有大小二。

\ProseParagraphBreak
还要注意，\(q\) 原来不可约，并不保证扩域后各 \(h_j=1\)。这种情况在正特征中会实际出现。取 \(k=\mathbb F_2(s)\)，其中 \(s\) 为不定元，即 \(k\) 是二元域上的有理函数域。多项式
\[
q(t)=t^2-s
\]
在 \(k[t]\) 中不可约。事实上，二次多项式若可约就有根；若 \(a(s)/b(s)\) 是根，其中 \(a,b\in\mathbb F_2[s]\)、\(b\ne0\)，则
\[
a(s)^2=s\cdot b(s)^2.
\]
两侧都非零，次数却分别为偶数 \(2\deg a\) 与奇数 \(1+2\deg b\)，矛盾。

\ProseParagraphBreak
现取另一个不定元 \(u\)，令 \(L=\mathbb F_2(u)\)，并通过 \(s\mapsto u^2\) 把 \(k\) 嵌入 \(L\)。这是嵌入，因为非零多项式代入超越元 \(u^2\) 后仍非零，故分式的分母不会变成零。在特征二中有
\[
q(t)=t^2-u^2=(t-u)^2.
\]
于是 \(q\) 在分裂域中只有一个二重根，是一个不可分多项式。矩阵
\[
A=C(q)=\begin{pmatrix}0&s\\1&0\end{pmatrix}
\]
在 \(k\) 上的初等因子是 \(q\)，即原不可约因子的一次幂；在 \(L\) 上却变成 \((t-u)^2\)，因而相似于 \(J_2(u)\)。也可直接检查 \((A-uI)^2=0\) 而 \(A-uI\ne0\)。这里 \(\mu_A=q\) 这个多项式在扩域后并未改变，只是同一个多项式由不可约因子 \(q\) 的一次幂，改写为线性因子 \(t-u\) 的二次幂。原指数 \(e=1\) 和根重数 \(h_1=2\) 相乘，才得到块大小 \(2\)；不可约因子的指数与分裂后线性因子的指数，只有在可分情形下才直接相同。
